% Vietnamese translation for OI Public Library
% Source-PDF: 数学 MATHEMATICS/容斥原理 - 张坤龙.pdf
\documentclass[11pt]{article}
\usepackage{oipl-vn}

\newcommand{\Oh}{\mathcal{O}}
\newcommand{\C}[2]{\binom{#1}{#2}}
\newcommand{\card}[1]{\left|#1\right|}

\title{Nguyên lý bao hàm - loại trừ}
\author{Zhang Kunlong}
\date{Toán tổ hợp ứng dụng, bài giảng 9}

\begin{document}
\maketitle

\origtitle{Inclusion-exclusion Principle}
\sourcepath{Mathematics / Inclusion-exclusion Principle - Zhang Kunlong}

\begin{abstract}
Bài giảng trình bày nguyên lý bao hàm - loại trừ từ công thức đếm hợp nhiều
tập, công thức Sylvester cho phần không có thuộc tính nào, đến dạng tổng quát
đếm phần tử có đúng \(m\) thuộc tính. Các ứng dụng gồm bội số, sàng
Eratosthenes, hàm Euler, hoán vị sai vị trí, hoán vị có vùng cấm, đa thức bàn
cờ xe, tổ hợp của đa tập, số Stirling loại hai, tổ hợp không kề nhau và bài
toán các cặp vợ chồng quanh bàn tròn.
\end{abstract}

\tableofcontents

\section{Nguyên lý bao hàm - loại trừ}

\subsection{Hai và ba tập}

Với hai tập hữu hạn \(A,B\), ta có
\[
  \card{A\cup B}=\card A+\card B-\card{A\cap B}.
\]
Ví dụ, số số nguyên dương không vượt quá \(20\) là bội của \(2\) hoặc \(3\)
là
\[
  \left\lfloor\frac{20}{2}\right\rfloor+
  \left\lfloor\frac{20}{3}\right\rfloor-
  \left\lfloor\frac{20}{6}\right\rfloor=10+6-3=13.
\]

Với ba tập:
\[
\begin{aligned}
  \card{A\cup B\cup C}
  ={}&\card A+\card B+\card C\\
  &-\card{A\cap B}-\card{A\cap C}-\card{B\cap C}\\
  &+\card{A\cap B\cap C}.
\end{aligned}
\]
Nếu một trường chỉ có ba môn toán, vật lý, hóa học, số học sinh học ba môn
lần lượt là \(170,130,120\), số học sinh học từng cặp lần lượt là \(45,20,22\),
và số học sinh học cả ba là \(3\), thì tổng số học sinh là
\[
  170+130+120-45-20-22+3=336.
\]

\subsection{Nhiều tập}

Với các tập hữu hạn \(A_1,A_2,\ldots,A_n\),
\[
  \card{\bigcup_{i=1}^n A_i}
  =
  \sum_i \card{A_i}
  -\sum_{i<j}\card{A_i\cap A_j}
  +\sum_{i<j<k}\card{A_i\cap A_j\cap A_k}
  -\cdots
  +(-1)^{n-1}\card{\bigcap_{i=1}^n A_i}.
\]
Công thức có thể chứng minh bằng quy nạp, hoặc bằng cách xét một phần tử nằm
trong đúng \(t\) tập. Nó được cộng vào
\[
  \C t1-\C t2+\C t3-\cdots+(-1)^{t-1}\C tt=1
\]
lần.

\subsection{Định lý De Morgan và công thức Sylvester}

Trong tập vũ trụ \(N\), ký hiệu \(\overline A=N\setminus A\). Với hai tập:
\[
  \overline{A\cup B}=\overline A\cap\overline B,\qquad
  \overline{A\cap B}=\overline A\cup\overline B.
\]
Tổng quát:
\[
  \overline{\bigcup_i A_i}=\bigcap_i\overline{A_i},\qquad
  \overline{\bigcap_i A_i}=\bigcup_i\overline{A_i}.
\]

Nếu \(A_i\) là tập các phần tử có thuộc tính \(i\), thì số phần tử không có
bất kỳ thuộc tính nào là
\[
  \card{\bigcap_{i=1}^n \overline{A_i}}
  =
  \card N-\sum_i\card{A_i}
  +\sum_{i<j}\card{A_i\cap A_j}
  -\cdots
  +(-1)^n\card{\bigcap_i A_i}.
\]
Công thức này thường được gọi là công thức Sylvester.

\section{Một số ứng dụng cơ bản}

\subsection{Hoán vị bị cấm chứa mẫu liên tiếp}

Xét các hoán vị của sáu chữ cái \(a,b,c,d,e,f\). Cần đếm số hoán vị không
chứa khối liên tiếp \(\texttt{ace}\) và không chứa khối liên tiếp
\(\texttt{df}\). Gọi \(A\) là tập hoán vị xem \(\texttt{ace}\) như một phần
tử, \(B\) là tập hoán vị xem \(\texttt{df}\) như một phần tử. Khi đó
\[
  \card N=6!,\quad \card A=4!,\quad \card B=5!,\quad
  \card{A\cap B}=3!,
\]
nên đáp án là
\[
  6!-(4!+5!)+3!=582.
\]

Với bốn ký tự \(x\), ba ký tự \(y\), hai ký tự \(z\), cần đếm các hoán vị đa
tập không xuất hiện các khối \(\texttt{xxxx}\), \(\texttt{yyy}\),
\(\texttt{zz}\). Ký hiệu
\[
  P(n;a,b,c)=\frac{n!}{a!\,b!\,c!}.
\]
Áp dụng bao hàm - loại trừ:
\[
\begin{aligned}
&P(9;4,3,2)
-\bigl(P(6;1,3,2)+P(7;4,1,2)+P(8;4,3,1)\bigr)\\
&+\bigl(P(4;1,1,2)+P(5;1,3,1)+P(6;4,1,1)\bigr)
-P(3;1,1,1)\\
&=1260-(60+105+280)+(12+20+30)-6=871.
\end{aligned}
\]

\subsection{Sàng Eratosthenes}

Cần đếm số nguyên tố không vượt quá \(120\). Mọi hợp số không vượt quá \(120\)
có ít nhất một ước nguyên tố trong \(\{2,3,5,7\}\). Gọi \(A_p\) là tập các số
không vượt quá \(120\) chia hết cho \(p\). Khi đó số không chia hết cho bất kỳ
phần tử nào trong \(\{2,3,5,7\}\) là
\[
  120-\sum_p\left\lfloor\frac{120}{p}\right\rfloor
  +\sum_{p<q}\left\lfloor\frac{120}{pq}\right\rfloor
  -\sum_{p<q<r}\left\lfloor\frac{120}{pqr}\right\rfloor
  +\left\lfloor\frac{120}{2\cdot3\cdot5\cdot7}\right\rfloor
  =27.
\]
Các số \(2,3,5,7\) là nguyên tố nhưng bị loại trong bước trên, còn \(1\)
không phải nguyên tố, nên số nguyên tố không vượt quá \(120\) là
\[
  27+4-1=30.
\]

\subsection{Hàm Euler}

Giả sử
\[
  n=p_1^{\alpha_1}p_2^{\alpha_2}\cdots p_k^{\alpha_k}.
\]
Hàm Euler \(\varphi(n)\) đếm số nguyên dương không vượt quá \(n\) và nguyên
tố cùng nhau với \(n\). Gọi \(A_i\) là tập các số từ \(1\) đến \(n\) chia hết
cho \(p_i\). Theo công thức Sylvester:
\[
\begin{aligned}
\varphi(n)
&=n-\sum_i\frac{n}{p_i}
 +\sum_{i<j}\frac{n}{p_ip_j}
 -\cdots
 +(-1)^k\frac{n}{p_1p_2\cdots p_k}\\
&=n\left(1-\frac1{p_1}\right)
    \left(1-\frac1{p_2}\right)\cdots
    \left(1-\frac1{p_k}\right).
\end{aligned}
\]

\subsection{Hoán vị sai vị trí}

Số hoán vị của \(1,2,\ldots,n\) sao cho không có \(i\) nào đứng ở vị trí
\(i\) là
\[
  D_n
  =n!\left(1-\frac1{1!}+\frac1{2!}-\frac1{3!}
  +\cdots+(-1)^n\frac1{n!}\right).
\]
Ở đây \(A_i\) là tập các hoán vị đặt \(i\) đúng vị trí. Nếu cố định \(k\) vị
trí thì còn \((n-k)!\) cách sắp các phần tử còn lại, nên công thức theo
bao hàm - loại trừ là
\[
  n!-\C n1(n-1)!+\C n2(n-2)!-\cdots+(-1)^n\C nn.
\]

\section{Hoán vị có vùng cấm và đa thức bàn cờ}

\subsection{Mô hình bàn cờ xe}

Một hoán vị \(P=P_1P_2\cdots P_n\) có thể xem như cách đặt \(n\) quân xe trên
bàn cờ \(n\times n\), mỗi hàng và mỗi cột đúng một quân. Ràng buộc
\(P_i\ne j\) trở thành một ô cấm tại hàng \(i\), cột \(j\).

Ví dụ với \(n=4\), các ràng buộc
\[
  P_1\ne3,\qquad P_2\ne1,4,\qquad
  P_3\ne2,4,\qquad P_4\ne2
\]
chính là bài toán đếm các cách đặt xe tránh các ô cấm tương ứng.

\subsection{Đa thức bàn cờ}

Với một tập ô \(B\) trên bàn cờ, đặt \(r_k(B)\) là số cách đặt \(k\) quân xe
trên các ô của \(B\) sao cho không có hai quân cùng hàng hoặc cùng cột. Đa
thức
\[
  R(B)=\sum_{k\ge0}r_k(B)x^k
\]
được gọi là đa thức bàn cờ. Vì bàn cờ hữu hạn nên \(r_k(B)=0\) với \(k\) đủ
lớn.

Nếu chọn một ô \(e\in B\), tách các cách đặt xe thành hai loại: có dùng ô
\(e\) và không dùng ô \(e\). Gọi \(B_e\) là bàn cờ còn lại sau khi xóa hàng
và cột chứa \(e\), còn \(B-e\) là bàn cờ sau khi chỉ xóa ô \(e\). Khi đó
\[
  r_k(B)=r_{k-1}(B_e)+r_k(B-e),
\]
và
\[
  R(B)=xR(B_e)+R(B-e).
\]
Đây là công thức truy hồi chuẩn để tính đa thức bàn cờ.

\subsection{Đếm cách tránh vùng cấm}

Giả sử vùng cấm của bàn cờ \(n\times n\) có các số \(r_i\) như trên. Số cách
đặt \(n\) quân xe, mỗi hàng mỗi cột đúng một quân, và không quân nào nằm trong
vùng cấm là
\[
  n!-r_1(n-1)!+r_2(n-2)!-\cdots+(-1)^n r_n.
\]
Chứng minh: gọi \(A_i\) là tập các hoán vị có quân xe thứ \(i\) rơi vào vùng
cấm ở hàng \(i\). Nếu đã chọn \(k\) quân nằm trong vùng cấm theo cách không
đụng hàng cột, \(n-k\) quân còn lại có \((n-k)!\) cách đặt. Áp dụng
bao hàm - loại trừ cho các sự kiện ``rơi vào vùng cấm'' sẽ cho công thức trên.

Trong bài phân công bốn công nhân \(G,L,W,Y\) cho bốn việc \(A,B,C,D\), nếu
\(G\) không làm được \(B\), \(L\) không làm được \(B,C\), \(W\) không làm được
\(C,D\), và \(Y\) không làm được \(D\), đa thức vùng cấm là
\[
  R(B)=1+6x+10x^2+4x^3.
\]
Do đó số phân công hợp lệ là
\[
  4!-6\cdot3!+10\cdot2!-4\cdot1!=4.
\]

Hoán vị sai vị trí cũng là một trường hợp vùng cấm: vùng cấm là đường chéo
chính, có đa thức
\[
  R(B)=(1+x)^n.
\]
Công thức vùng cấm vì thế quay lại đúng công thức sai vị trí ở trên.

\section{Nguyên lý bao hàm - loại trừ tổng quát}

\subsection{Ký hiệu}

Cho tập \(N\) và các thuộc tính \(A_1,\ldots,A_n\). Đặt
\[
  \alpha(0)=\card N,
\]
và với \(m\ge1\),
\[
  \alpha(m)=
  \sum_{1\le i_1<i_2<\cdots<i_m\le n}
  \card{A_{i_1}\cap A_{i_2}\cap\cdots\cap A_{i_m}}.
\]
Một phần tử có đúng \(m+k\) thuộc tính sẽ được tính trong \(\alpha(m)\) đúng
\(\C{m+k}{m}\) lần.

Gọi \(\beta(m)\) là số phần tử của \(N\) có đúng \(m\) thuộc tính. Khi đó
\[
  \beta(m)
  =
  \alpha(m)-\C{m+1}{m}\alpha(m+1)
  +\C{m+2}{m}\alpha(m+2)
  -\cdots
  +(-1)^{n-m}\C nm\alpha(n).
\]
Trường hợp \(m=0\) chính là công thức Sylvester:
\[
  \beta(0)=\alpha(0)-\alpha(1)+\alpha(2)-\cdots+(-1)^n\alpha(n).
\]

\subsection{Ví dụ về giáo viên}

Một trường có \(12\) giáo viên. Có \(8\) người dạy toán, \(6\) người dạy vật
lý, \(5\) người dạy hóa; dạy toán-vật lý \(5\), toán-hóa \(4\), vật lý-hóa
\(3\); dạy cả ba môn \(3\). Khi đó
\[
  \alpha(0)=12,\quad \alpha(1)=8+6+5=19,\quad
  \alpha(2)=5+4+3=12,\quad \alpha(3)=3.
\]
Số giáo viên dạy môn khác, không dạy môn nào trong ba môn trên:
\[
  \beta(0)=12-19+12-3=2.
\]
Số giáo viên chỉ dạy một trong ba môn:
\[
  \beta(1)=\alpha(1)-2\alpha(2)+3\alpha(3)=19-24+9=4.
\]
Số giáo viên chỉ dạy đúng hai môn:
\[
  \beta(2)=\alpha(2)-3\alpha(3)=12-9=3.
\]

\subsection{Đường đi qua đúng hai đoạn bắt buộc}

Với các đường đi lưới từ \(O\) đến \(P\), giả sử có bốn đoạn đặc biệt
\(AB,CD,EF,GH\). Gọi \(A_i\) là tập các đường đi qua đoạn thứ \(i\). Bài
giảng tính được
\[
  \alpha(2)=861,\qquad \alpha(3)=150,\qquad \alpha(4)=0.
\]
Số đường đi qua đúng hai trong bốn đoạn là
\[
  \beta(2)=\alpha(2)-3\alpha(3)+6\alpha(4)=861-450=411.
\]

\subsection{Một đồng nhất thức tổ hợp}

Công thức tổng quát cũng chứng minh các đồng nhất thức nhị thức. Ví dụ, số
cách chọn \(k\) số từ \(\{1,2,\ldots,n\}\) và bắt buộc chứa
\(\{1,2,\ldots,m\}\) là \(\C{n-m}{k-m}\). Nếu \(A_i\) là tính chất ``không
chọn số \(i\)'', với \(1\le i\le m\), áp dụng bao hàm - loại trừ cho các tính
chất bị thiếu sẽ cho khai triển luân phiên tương ứng của \(\C{n-m}{k-m}\).

\section{Phương trình nguyên và đa tập}

\subsection{Số nghiệm không âm}

Số nghiệm nguyên không âm của
\[
  x_1+x_2+x_3=15
\]
là số cách đặt \(15\) quả bóng giống nhau vào \(3\) hộp phân biệt:
\[
  \C{15+3-1}{3-1}=\C{17}{2}=136.
\]
Nếu thêm ràng buộc \(x_1\ge6\), đặt \(y_1=x_1-6\), \(y_2=x_2\), \(y_3=x_3\),
ta được
\[
  y_1+y_2+y_3=9,
\]
nên số nghiệm là \(\C{11}{2}=55\).

Nếu ràng buộc là \(0\le x_1\le5\), lấy toàn bộ nghiệm không âm rồi trừ các
nghiệm \(x_1\ge6\):
\[
  \C{17}{2}-\C{11}{2}=136-55=81.
\]

\subsection{Tổ hợp của đa tập}

Với đa tập \(\{5\cdot x,6\cdot y,7\cdot z\}\), số cách chọn \(15\) phần tử là
số nghiệm của
\[
  x_1+x_2+x_3=15,\qquad
  0\le x_1\le5,\quad 0\le x_2\le6,\quad 0\le x_3\le7.
\]
Lấy \(N\) là tập mọi nghiệm không âm, \(A_1\) là tập nghiệm có \(x_1\ge6\),
\(A_2\) có \(x_2\ge7\), \(A_3\) có \(x_3\ge8\). Bài giảng tính
\[
  \alpha(0)=136,\qquad \alpha(1)=136,\qquad
  \alpha(2)=10,\qquad \alpha(3)=0,
\]
nên
\[
  \beta(0)=136-136+10=10.
\]

\section{Số Stirling loại hai}

Đặt \(S(n,m)\) là số cách chia \(n\) phần tử phân biệt thành \(m\) khối không
rỗng. Nếu trước hết bỏ \(n\) quả bóng phân biệt vào \(m\) hộp phân biệt, số
cách không để hộp nào rỗng là \(m!\,S(n,m)\).

Gọi \(A_i\) là sự kiện hộp \(i\) rỗng. Khi đó
\[
  \alpha(k)=\C mk(m-k)^n.
\]
Do đó
\[
  m!\,S(n,m)
  =
  \sum_{k=0}^m (-1)^k\C mk(m-k)^n,
\]
hay
\[
  S(n,m)=\frac1{m!}\sum_{k=0}^m(-1)^k\C mk(m-k)^n.
\]

\section{Các mở rộng khác}

\subsection{Sai vị trí từng phần}

Gọi \(D(n,r,k)\) là số cách lấy \(r\) phần tử từ \(N=\{1,2,\ldots,n\}\) rồi
sắp thành dãy \(a_1a_2\cdots a_r\), trong đó đúng \(k\) vị trí thỏa
\(a_i=i\). Với \(n\ge r\ge k\), bài giảng đưa ra công thức
\[
  D(n,r,k)
  =
  \frac{\C rk}{(n-r)!}
  \sum_{i=0}^{r-k}(-1)^i\C{r-k}{i}(n-k-i)!.
\]
Ý tưởng chứng minh là dùng \(A_i\) cho thuộc tính \(a_i=i\), tính
\(\alpha(t)\), rồi lấy \(\beta(k)\) bằng bao hàm - loại trừ tổng quát.

\subsection{Tổ hợp không kề nhau}

Số cách chọn \(r\) số từ \(\{1,2,\ldots,n\}\) sao cho không có hai số kề nhau
là
\[
  \C{n-r+1}{r}.
\]
Song ánh chuẩn là nếu chọn
\[
  b_1<b_2<\cdots<b_r
\]
không kề nhau trong \(\{1,\ldots,n\}\), thì ánh xạ sang
\[
  b_1,\ b_2-1,\ b_3-2,\ldots,\ b_r-(r-1)
\]
là một tổ hợp thường của \(r\) phần tử trong \(\{1,\ldots,n-r+1\}\).

Nếu \(n\) đỉnh nằm trên một vòng tròn, số cách chọn \(k\) đỉnh không kề nhau
là
\[
  \frac{n}{k}\C{n-k-1}{k-1}.
\]
Chứng minh: cố định một đỉnh \(A\) được chọn. Sau khi loại hai đỉnh kề \(A\),
còn \(n-3\) đỉnh trên một đoạn thẳng, cần chọn \(k-1\) đỉnh không kề nhau, có
\(\C{n-k-1}{k-1}\) cách. Vì mỗi tập chọn có \(k\) đỉnh và có thể được đếm từ
mỗi đỉnh được chọn, nhân \(n\) rồi chia \(k\).

\subsection{Bài toán n cặp vợ chồng}

Có \(n\) cặp vợ chồng ngồi quanh bàn tròn, yêu cầu nam nữ xen kẽ và vợ chồng
không ngồi cạnh nhau. Đầu tiên xếp \(n\) người vợ quanh bàn; giữa hai người
vợ có một chỗ cho chồng. Đánh số các chỗ trống theo chiều kim đồng hồ.

Với mỗi cặp \(i\), có hai thuộc tính xấu: người chồng \(i\) ngồi bên phải vợ
\(i\), hoặc ngồi bên trái vợ \(i\). Hai thuộc tính xấu của cùng một cặp không
thể đồng thời xảy ra. Nếu \(1\le k\le n\), số cách chọn \(k\) thuộc tính xấu
không xung đột trên vòng là
\[
  \frac{2n}{k}\C{2n-k-1}{k-1}.
\]
Sau khi cố định \(k\) thuộc tính xấu, còn \((n-k)!\) cách đặt các người chồng
còn lại. Vì vậy
\[
  \alpha(k)=\frac{2n}{k}\C{2n-k-1}{k-1}(n-k)!.
\]
Theo bao hàm - loại trừ, số cách không có cặp vợ chồng nào ngồi cạnh nhau là
\[
  \beta(0)=
  \sum_{k=0}^{n}
  (-1)^k\frac{2n}{2n-k}\C{2n-k}{k}(n-k)!,
\]
với hạng \(k=0\) được hiểu là \(n!\). Cuối cùng, nếu tính cả việc xếp các
người vợ quanh bàn tròn, nhân thêm \((n-1)!\).

\section*{Kết luận}

Bao hàm - loại trừ không chỉ là công thức đếm hợp của vài tập. Nó là một cách
tổ chức lại việc đếm: bắt đầu từ một không gian dễ đếm, rồi trừ và cộng lại
các trạng thái vi phạm theo số thuộc tính. Dạng tổng quát cho phép đếm đúng số
thuộc tính, còn mô hình bàn cờ xe biến các bài hoán vị có vùng cấm thành bài
toán tính đa thức bàn cờ.

\end{document}
